\documentclass[11pt]{article}

% ---------------------------------------------------------------
% Extending the Actor Model with Synchronous Delivery:
% fast_send from CPU Threads to FPGA Pipelines
% ---------------------------------------------------------------

\usepackage[utf8]{inputenc}
\usepackage[T1]{fontenc}
\usepackage{lmodern}
\usepackage{amsmath,amssymb,amsthm}
\usepackage{mathpartir}
\usepackage{booktabs}
\usepackage[margin=1in]{geometry}
\usepackage{hyperref}
\usepackage{xcolor}
\usepackage{enumitem}
\usepackage{microtype}
\usepackage{caption}
\usepackage[capitalise,noabbrev]{cleveref}

\hypersetup{colorlinks=true,linkcolor=blue!50!black,citecolor=blue!50!black,urlcolor=blue!50!black}

\newtheorem{theorem}{Theorem}
\newtheorem{lemma}[theorem]{Lemma}
\newtheorem{proposition}[theorem]{Proposition}
\newtheorem{corollary}[theorem]{Corollary}
\theoremstyle{definition}
\newtheorem{definition}[theorem]{Definition}
\newtheorem{assumption}[theorem]{Assumption}
\theoremstyle{remark}
\newtheorem{remark}[theorem]{Remark}

\newcommand{\fs}{\texttt{fast\_send}}
\newcommand{\snd}{\texttt{send}}
\newcommand{\kh}{kaspar-hft}

% semantic entities
\newcommand{\Act}[1]{\mathcal{A}\langle #1\rangle}
\newcommand{\Thr}[1]{\mathcal{T}\langle #1\rangle}
\newcommand{\Frm}[1]{\mathcal{F}\langle #1\rangle}
\newcommand{\Msg}[1]{\mathcal{M}\langle #1\rangle}
\newcommand{\Obj}[1]{\mathcal{O}\langle #1\rangle}
\newcommand{\ctx}[1]{\mathcal{E}[#1]}
\newcommand{\turn}{\mathsf{turn}}
\newcommand{\ret}[1]{\mathsf{ret}(#1)}
\newcommand{\wait}[1]{\mathsf{wait}(#1)}
\newcommand{\fsend}{\mathsf{fastsend}}
\newcommand{\asend}{\mathsf{send}}
\newcommand{\busy}{\mathsf{busy}}
\newcommand{\body}{\mathsf{body}}
\newcommand{\pass}{\mathsf{pass}}
\newcommand{\red}{\longrightarrow}

\title{Extending the Actor Model with Synchronous Delivery:\\
\fs{} from CPU Threads to FPGA Pipelines}
\author{Vincent Mayeski\\M2 Tech, Montreal\\\texttt{mayeski@gmail.com}}
\date{October 2026}

\begin{document}
\maketitle

\begin{abstract}
The actor model has one communication primitive, asynchronous send. Practical actor
systems add synchronous request/reply on top of it, as a library pattern (Erlang's
\texttt{gen\_server:call}, Akka's \texttt{ask}) or as a call restricted to a single
scheduling unit (ABS cogs, E vats). We add a second primitive to the model itself:
\fs{}, in which the receiver's turn runs now, the caller waits for it, and the reply
comes back as a value. We give \fs{} an operational semantics by extending the
active-object calculus of De Koster and De Meuter with explicit execution resources, and
we compare it with their isolated turn principle, which has a safety half (a turn is
isolated and its result is determined) and a liveness half (a turn never blocks). \fs{}
keeps isolation: no two turns of an actor overlap and no turn touches another actor's
state. It weakens determinism to determinism relative to the replies a turn receives.
It gives up unconditional liveness: we prove that a turn fails to complete exactly when
its thread reaches a cycle in a wait-for graph, and that with asynchronous send alone
the graph is empty, which recovers the original result. A per-thread call-chain check
sees only the cycles that stay on one thread. We show that \fs{} is not reducible to
asynchronous send followed by a blocking wait once execution resources are shared, which
is where the classic encodings of synchronous into asynchronous communication do not
apply. For debugging and backtesting, we prove that running every actor --- whatever its
language or substrate --- from one global queue is sound: each execution of that
simulator is an execution of the concurrent system, \fs{} included, so a failure seen
in simulation is real, while some failures of the concurrent system, cross-thread
deadlock among them, cannot appear in it. Finally we map \snd{} and \fs{} onto CPU
threads, an FPGA actor runtime, a Rust runtime and, as a design, a GPU runtime: on each
substrate the two primitives correspond to two distinct mechanisms, a queue and an
inline or blocking call. We report their measured cost on the CPU runtime.
\end{abstract}

% ---------------------------------------------------------------
\section{Introduction}
% ---------------------------------------------------------------

An actor owns private state and communicates only by messages. In the classic
formulations \cite{hewitt1973,agha1986,agha1997} there is exactly one way to
communicate: an asynchronous send, which places a message in the receiver's mailbox and
returns at once. The receiver processes its messages one at a time, each in a
\emph{turn}. Because a turn never waits on anything, turns cannot deadlock, and because
no state is shared, turns cannot race. De Koster and De Meuter name this the
\emph{isolated turn principle}, give an operational semantics for four families of actor
languages, and prove the principle for each \cite{dekoster2024}.

Real actor systems nonetheless need request/reply, and they obtain it in one of two ways.
Either the sender sends a request and then blocks until a reply message arrives --- the
receiver still runs on its own thread (ABCL/1's \emph{now}-type messages
\cite{yonezawa1986}, ABS calls between cogs \cite{johnsen2012abs}, Erlang's
\texttt{gen\_server:call}, Akka's \texttt{ask}) --- or a synchronous call is allowed only
between objects that already share one scheduling unit (ABS calls inside a cog, E's
immediate calls inside a vat \cite{miller2005}).

This paper studies a third form, which we call \fs{} and which the \kh{} framework
provides \cite{mayeski2026fastsend}: the receiver's turn is run \emph{now}, by an
execution resource that the caller is willing to wait on --- on a CPU, the caller's own
thread --- under the receiver's exclusive access, and its result is returned to the
caller as a value. It is neither a blocking wait on another thread's work nor a call
confined to one scheduling unit. We argue that it belongs in the model rather than in a
library, for two reasons. First, its meaning can be stated precisely and its effect on
the model's guarantees characterised exactly (Sections~\ref{sec:calculus}
and~\ref{sec:properties}). Second, it maps onto a distinct physical mechanism on every
substrate we have built, just as asynchronous send does (Section~\ref{sec:substrates}).

\paragraph{Contributions.}
\begin{enumerate}[leftmargin=1.6em]
  \item A thread-aware extension of the active-object calculus of
  \cite{dekoster2024} with a \fs{} rule, in two variants: the callee's turn runs on the
  caller's resource (\textsc{Fast-Local}) or on the callee's own resource while the
  caller waits (\textsc{Fast-Remote}) (Section~\ref{sec:calculus}).
  \item Its properties, stated against the isolated turn principle
  (Section~\ref{sec:properties}): isolation is kept (\cref{thm:isolation});
  determinism holds relative to replies (\cref{thm:determinism}); liveness holds
  exactly when no thread reaches a cycle of waits (\cref{thm:liveness}), and the
  original liveness result is the special case without \fs{}
  (\cref{cor:asynclive}); a per-thread call-chain check detects exactly the cycles
  confined to one thread (\cref{prop:localcycles}); \fs{} overtakes queued messages
  (\cref{prop:order}); and \fs{} is not send-then-wait when resources are shared
  (\cref{prop:notsendwait}).
  \item A soundness theorem for debugging and backtesting from a single global queue,
  across languages and substrates (\cref{thm:gqsound}), its limits
  (\cref{prop:gqincomplete}), and an offline check for the deadlocks it cannot
  exhibit (\cref{prop:callgraph}) (Section~\ref{sec:debug}).
  \item A mapping of both primitives onto CPU threads, an FPGA actor runtime, a Rust
  runtime and a GPU runtime design, and between them (Section~\ref{sec:substrates}),
  with measurements of the CPU runtime (Section~\ref{sec:eval}).
\end{enumerate}

% ---------------------------------------------------------------
\section{Background: turns and the isolated turn principle}
\label{sec:background}
% ---------------------------------------------------------------

De Koster and De Meuter classify actor languages into four families --- classic
actors, active objects, processes and communicating event loops --- and give each a
small-step operational semantics \cite{dekoster2024}. We build on their active-object
calculus, because it matches the programming model of the runtimes studied here: an
actor has a fixed interface, a FIFO mailbox and imperative private state.

In that calculus a configuration is a set of actors $\Act{\iota_a, \mu, e, O, o}$: an
identifier, a mailbox $\mu$ of messages $\Msg{msg, \overline{v}}$, the expression $e$
the actor is evaluating, a heap $O$ of passive objects it owns, and an active object $o$
that fixes its interface. Actor-local rules (\textsc{Let}, \textsc{New},
\textsc{Invoke}, \textsc{Field-Access}, \textsc{Field-Update}, \textsc{Self-Send},
\textsc{Receive}) reduce one actor; actor-global rules (\textsc{Send}, \textsc{Spawn})
involve two. \textsc{Send} appends a message to the receiver's mailbox after deep-copying
the passive objects it references (the function $\pass$), so no state is shared.
\textsc{Receive} starts a turn: when an actor's expression is a value, the head of its
mailbox is removed and the matching method body becomes its expression. A \emph{turn} is
the sequence of reductions of one actor between two applications of \textsc{Receive}.

The isolated turn principle then consists of two theorems \cite{dekoster2024}.
\emph{Safety}: the value a turn reduces to is determined up to $\alpha$-equivalence,
whatever the interleaving with other actors, and a turn is not interleaved with another
turn of the same actor. \emph{Liveness}: a turn never blocks --- every reducible
expression of a turn can take a step --- so turns are free of deadlock. The calculi
contain no blocking operation; this is what makes liveness hold.

Two features of this setting matter below. The calculus has no notion of the thread or
processor that performs a reduction: each actor reduces independently, as if it had its
own. And every communication is a \textsc{Send}. A synchronous call needs the first to
be made explicit and adds a second kind of communication.

% ---------------------------------------------------------------
\section{Synchronous communication in actor systems}
\label{sec:related-sync}
% ---------------------------------------------------------------

Three reasons are usually given for leaving synchronous calls out of the actor model.
\emph{Order}: in Hewitt's model messages are guaranteed to arrive but in no particular
order, so ``handle this now, before what I do next'' has no meaning. \emph{Deadlock}: an
actor that handles one message at a time and waits inside a turn can wait on itself,
directly or through others. \emph{Distribution}: across machines, a synchronous call can
wait forever on a failed peer.

Where systems provide synchronous communication anyway, they do so in one of two ways.

\paragraph{Wait on the receiver's thread.} ABCL/1 distinguishes \emph{past}
(asynchronous), \emph{now} (send and wait for the reply) and \emph{future} messages
\cite{yonezawa1986}. In ABS \cite{johnsen2012abs}, a synchronous call $v = o.m(e)$
between objects in different cogs abbreviates $x = o!m(e);\ v = x.\mathtt{get}$: an
asynchronous call followed by a blocking \texttt{get}, which ``blocks all execution in the
object until the return value is available''. Erlang's \texttt{gen\_server:call} and
Akka's \texttt{ask} are library forms of the same pattern, usually with a timeout. In all
of these the receiver's turn runs on the receiver's own thread; the caller's thread
waits.

\paragraph{Call within one scheduling unit.} In ABS, a synchronous call between objects
in the same cog has ``the reentrant semantics known from, e.g., Java threads''
\cite{johnsen2012abs}: control passes to the callee on the cog's single processor. In E,
immediate calls are allowed only on near references, to objects in the same vat; between
vats only eventual sends exist \cite{miller2005}. These calls run on the caller's
processor, but they never cross a scheduling unit, so they need no mutual exclusion; and
ABS allows re-entry into an object that is already executing.

\paragraph{Expressiveness.} Honda and Tokoro \cite{hondatokoro1991} and Boudol
\cite{boudol1992} showed that synchronous communication can be encoded in the
asynchronous $\pi$-calculus. \fs{} therefore adds nothing to what can be computed. The
encodings assume, however, that every process makes progress independently --- that
execution resources are unbounded. \Cref{prop:notsendwait} shows that once resources are
shared, as in an actor group or a cog, send-then-wait deadlocks where \fs{} completes. The
difference between the primitives is a difference in who executes a turn and when.

\paragraph{Where \fs{} sits.} \fs{} runs the receiver's turn on a resource the caller is
willing to wait on --- on a CPU, the caller's own thread --- across scheduling units,
under the receiver's mutual exclusion, and without re-entry: an actor already in a turn
cannot be entered again, and a call that tries waits, possibly forever. It combines the
``run it here'' of the second group with the ``any actor'' of the first, and pays for it
with the deadlock case that the first group has and the second avoids. The rest of the
paper makes this precise.

% ---------------------------------------------------------------
\section{A thread-aware calculus with \fs{}}
\label{sec:calculus}
% ---------------------------------------------------------------

We extend the active-object calculus of \cite{dekoster2024} in three ways: execution
resources become explicit, an actor's current expression moves from the actor into a
\emph{frame} on a resource's stack, and a \fs{} expression is added. We call the
resources \emph{threads}; on an FPGA they are hardware processes
(Section~\ref{sec:substrates}).

\subsection{Semantic entities and syntax}

\begin{align*}
K &\in \text{Configuration} ::= \langle A, T\rangle\\
a &\in A \subseteq \text{Actor} ::= \Act{\iota_a, \mu, O, o, \theta}\\
t &\in T \subseteq \text{Thread} ::= \Thr{\theta, S}\qquad S ::= \epsilon \mid S\cdot f\\
f &\in \text{Frame} ::= \Frm{\iota_a, e, \kappa}\qquad \kappa ::= \turn \mid \ret{\rho}
\end{align*}

An actor has an identifier, a FIFO mailbox $\mu$, a heap $O$, an active object $o$, and
a \emph{home thread} $\theta$. Several actors may share a home thread; that is an actor
group. A thread has an identifier and a stack $S$ of frames, top on the right. A frame
records which actor it belongs to, the expression being evaluated, and why it exists:
$\turn$ for a turn started from the mailbox, or $\ret{\rho}$ for a turn started by the
\fs{} call $\rho$, whose value must be returned. Call identifiers $\rho$ are fresh.

The expression syntax of the active-object calculus is extended with one form, and the
run-time syntax with one more:
\[
e ::= \dots \mid \fsend(e, msg, \overline{e})\qquad\qquad
e ::= \dots \mid \wait{\rho}\ \ \text{(run time)}
\]
Evaluation contexts $\mathcal{E}$ extend the originals left to right through the
arguments of $\fsend$. No rule reduces $\wait{\rho}$: a frame whose expression has the
form $\ctx{\wait{\rho}}$ is \emph{waiting}.

\begin{definition}[Busy]
Actor $\iota_a$ is \emph{busy} in $K$, written $\busy_K(\iota_a)$, if some thread of $K$
has a frame of $\iota_a$ on its stack.
\end{definition}

For an actor with active object $o = \Obj{\iota_o, cls, \overline{v_f}}$ and a message
$\Msg{msg, \overline{v}}$, let $\body(a, msg, \overline{v})$ be the body of method $msg$
of $cls$ with fields, parameters and $\mathsf{self}$ substituted, exactly as in
\textsc{Receive} and \textsc{Invoke} of \cite{dekoster2024}.

\subsection{Reduction rules}

The actor-local rules of \cite{dekoster2024} apply unchanged to the \emph{top} frame of
any thread, reading and writing the heap of the frame's actor. \textsc{Send} and
\textsc{Self-Send} likewise apply to a top frame and append to the receiver's mailbox.
\textsc{Spawn} creates an actor with an empty mailbox and a home thread, either an
existing one or a fresh $\Thr{\theta', \epsilon}$. \textsc{Receive} and the end of a turn
become thread rules, and three rules are new (\cref{fig:rules}).

\begin{figure}[t]
\small
\begin{mathpar}
\inferrule*[right=Receive]
  {t = \Thr{\theta, \epsilon} \\ a = \Act{\iota_a, \Msg{msg,\overline{v}}\cdot\mu, O, o, \theta} \\ \neg\busy_K(\iota_a)}
  {K \red K[t \mapsto \Thr{\theta, \Frm{\iota_a, \body(a,msg,\overline{v}), \turn}}][a \mapsto \Act{\iota_a, \mu, O, o, \theta}]}

\inferrule*[right=Turn-End]
  {t = \Thr{\theta, S\cdot\Frm{\iota_a, v, \turn}}}
  {K \red K[t \mapsto \Thr{\theta, S}]}

\inferrule*[right=Fast-Local]
  {t = \Thr{\theta, S\cdot\Frm{\iota_a, \ctx{\fsend(\iota_b, msg, \overline{v})}, \kappa}} \\
   \neg\busy_K(\iota_b) \\ \overline{v}' = \pass(\overline{v}) \\ \rho\ \text{fresh}}
  {K \red K[t \mapsto \Thr{\theta, S\cdot\Frm{\iota_a, \ctx{\wait{\rho}}, \kappa}\cdot\Frm{\iota_b, \body(b,msg,\overline{v}'), \ret{\rho}}}]}

\inferrule*[right=Fast-Remote]
  {t = \Thr{\theta, S\cdot\Frm{\iota_a, \ctx{\fsend(\iota_b, msg, \overline{v})}, \kappa}} \\
   b = \Act{\iota_b, \mu_b, O_b, o_b, \theta_b} \\ \theta_b \neq \theta \\ t_b = \Thr{\theta_b, \epsilon} \\\\
   \neg\busy_K(\iota_b) \\ \overline{v}' = \pass(\overline{v}) \\ \rho\ \text{fresh}}
  {K \red K[t \mapsto \Thr{\theta, S\cdot\Frm{\iota_a, \ctx{\wait{\rho}}, \kappa}}][t_b \mapsto \Thr{\theta_b, \Frm{\iota_b, \body(b,msg,\overline{v}'), \ret{\rho}}}]}

\inferrule*[right=Fast-Return]
  {t' = \Thr{\theta', S'\cdot\Frm{\iota_b, v, \ret{\rho}}} \\
   t = \Thr{\theta, S\cdot\Frm{\iota_a, \ctx{\wait{\rho}}, \kappa}\cdot S''} \\ v' = \pass(v)}
  {K \red K[t' \mapsto \Thr{\theta', S'}][t \mapsto \Thr{\theta, S\cdot\Frm{\iota_a, \ctx{v'}, \kappa}\cdot S''}]}
\end{mathpar}
\caption{Thread rules. $K[t\mapsto t']$ replaces thread $t$ by $t'$ (and likewise for
actors). In \textsc{Fast-Return}, $t$ and $t'$ may be the same thread, in which case
$S''$ is empty and the returned frame was directly above the waiting one.}
\label{fig:rules}
\end{figure}

\begin{itemize}[leftmargin=1.4em]
  \item \textsc{Receive}: an idle thread starts a turn of one of its idle actors whose
  mailbox is not empty. With several actors on one thread, the choice among them is
  free.
  \item \textsc{Turn-End}: a turn started from the mailbox that has reduced to a value
  is removed. Its value is discarded; a reply to a \snd{} is an explicit \snd{}.
  \item \textsc{Fast-Local}: the caller's frame becomes waiting, and the callee's turn is
  pushed on the \emph{caller's} stack. This is \fs{} on one CPU thread, or inside one
  FPGA process.
  \item \textsc{Fast-Remote}: the caller's frame becomes waiting, and the callee's turn
  starts on the callee's own, idle, home thread. This is \fs{} to an actor on another
  resource that the caller waits on: a CPU thread calling into an FPGA process, an FPGA
  process calling a CPU actor, or a call to another machine.
  \item \textsc{Fast-Return}: a callee turn that has reduced to a value is removed and its
  value, copied, replaces the $\wait{\rho}$ of its caller.
\end{itemize}

Both \fs{} rules require the callee to be idle. When it is busy, neither applies: the
calling frame is \emph{blocked}. The callee's mailbox is not consulted. A system may use
\textsc{Fast-Local} for some pairs of actors and \textsc{Fast-Remote} for others; every
result below holds for any such mixture.

\begin{remark}[Receiver transparency]
A method body cannot tell how its turn started: \textsc{Receive}, \textsc{Fast-Local} and
\textsc{Fast-Remote} all install the same $\body(b, msg, \overline{v})$. This is the
calculus counterpart of receiver transparency in \kh{} \cite{mayeski2026fastsend}. It is
a property of the interface a turn is given, not a claim about what code in a real
language can observe about its execution context.
\end{remark}

% ---------------------------------------------------------------
\section{Properties}
\label{sec:properties}
% ---------------------------------------------------------------

Throughout, $K_0 \red^* K$ ranges over configurations reachable from an initial
configuration with one actor, as in \cite{dekoster2024}, whose single thread holds no
frames.

\subsection{Isolation}

\begin{lemma}[One frame per actor]
\label[lemma]{lem:oneframe}
In every reachable configuration, each actor has at most one frame, and every frame reads
and writes only the heap of its own actor.
\end{lemma}
\begin{proof}
Initially there are no frames. Only \textsc{Receive}, \textsc{Fast-Local} and
\textsc{Fast-Remote} create frames, and each creates a frame for an actor that is not
busy, that is, has no frame. \textsc{Turn-End} and \textsc{Fast-Return} only remove
frames; the other rules rewrite the expression of an existing top frame. The heap
discipline is that of the active-object calculus: field access and update apply to the
heap of the frame's actor, and values crossing actors (\textsc{Send}, the two \fs{} rules,
\textsc{Fast-Return}) are copied by $\pass$.
\end{proof}

\begin{definition}[Turn]
A turn of $\iota_a$ is the sequence of reductions from the creation of a frame of
$\iota_a$ (by \textsc{Receive}, \textsc{Fast-Local} or \textsc{Fast-Remote}) to its
removal (by \textsc{Turn-End} or \textsc{Fast-Return}).
\end{definition}

\begin{theorem}[Isolation]
\label{thm:isolation}
No two turns of the same actor overlap, and during a turn of $\iota_a$ the heap of
$\iota_a$ is changed only by reductions of that turn's frame.
\end{theorem}
\begin{proof}
Immediate from \cref{lem:oneframe}: a second turn would need a second frame.
\end{proof}

This is the safety half of the isolated turn principle, and it holds with \fs{}. A turn
that makes a \fs{} is suspended while the callee runs, but no other turn of the same
actor can run in the meantime, and the callee cannot touch the caller's state.

\subsection{Determinism relative to replies}

In \cite{dekoster2024} a turn's value is determined, because a turn receives no input
after it starts. With \fs{} it does: each \textsc{Fast-Return} substitutes a value that
depends on the callee's state, which other turns may have changed before the call.

\begin{theorem}[Determinism relative to replies]
\label{thm:determinism}
Consider two turns of the same actor that start from $\alpha$-equivalent heaps and
expressions. If the values substituted into them by \textsc{Fast-Return}, in order, are
$\alpha$-equivalent, then the two turns end with $\alpha$-equivalent values and heaps and
issue $\alpha$-equivalent \snd{} and \fs{} requests in the same order.
\end{theorem}
\begin{proof}
As the safety proof of \cite{dekoster2024}, by induction on the reductions of the turn's
frame: evaluation contexts fix the order of reduction, and every actor-local rule
produces $\alpha$-equivalent results from $\alpha$-equivalent inputs. The new cases are
the two \fs{} rules, which rewrite the frame to $\ctx{\wait{\rho}}$ and perform no
reduction on it until \textsc{Fast-Return}, whose substituted value is
$\alpha$-equivalent by hypothesis. Reductions of other frames do not touch this turn's
heap (\cref{lem:oneframe}).
\end{proof}

A turn that makes no \fs{} satisfies the original safety theorem unchanged. A turn that
does is deterministic as a function of its replies, as it would be if the replies were
messages; what \fs{} removes is the possibility of any other message of the same actor
being handled between the request and the reply.

\subsection{Liveness}

\begin{definition}[Blocked frame, wait-for graph]
A top frame $\Frm{\iota_a, \ctx{\fsend(\iota_b, msg, \overline{v})}, \kappa}$ on thread
$\theta$ is \emph{blocked} in $K$ if no \fs{} rule applies to it. The \emph{wait-for
graph} $W_K$ has the threads of $K$ as nodes and an edge $\theta \to \theta'$ when either
\begin{enumerate}[label=(\roman*),leftmargin=2em]
  \item the top frame of $\theta$ is blocked on $\iota_b$ and $\theta'$ holds the frame
  of $\iota_b$; or it is blocked by \textsc{Fast-Remote} because $\iota_b$'s home thread
  $\theta'$ is not idle; or
  \item $\theta$ holds a waiting frame $\ctx{\wait{\rho}}$ whose callee frame $\ret{\rho}$
  is on $\theta' \neq \theta$.
\end{enumerate}
\end{definition}

Every edge leaves a thread that cannot progress until something on the target thread is
removed. A thread whose top frame is waiting has an edge of kind (ii) unless the callee is
above it on the same stack, in which case the thread is not stuck: its top frame is the
callee's.

\begin{assumption}[Terminating handlers, fair scheduling]
\label[assumption]{ass:fair}
Every method body terminates when its \fs{} calls return, and every thread whose top
frame can be reduced is eventually reduced.
\end{assumption}

\begin{theorem}[Liveness exactly when waits are acyclic]
\label{thm:liveness}
Under \cref{ass:fair}, a turn in a reachable configuration $K$ fails to complete in some
fair continuation of $K$ if and only if, in some configuration reachable from $K$, its
thread has a path in the wait-for graph to a cycle.
\end{theorem}
\begin{proof}
($\Leftarrow$) Take the cycle $\theta_1 \to \dots \to \theta_n \to \theta_1$. Each
$\theta_i$ waits for a frame on $\theta_{i+1}$ to be removed. That frame is below the top
of $\theta_{i+1}$ or is its top; in either case its removal needs $\theta_{i+1}$ to
progress, which needs a frame on $\theta_{i+2}$ removed, and so on around the cycle. No
rule removes a frame of a thread whose top frame is blocked or waiting, so none of the
$\theta_i$ ever progresses, and a thread with a path to the cycle waits on it forever.

($\Rightarrow$) Suppose the thread $\theta$ of a turn never completes it, and that no
configuration reachable from $K$ has a path from $\theta$ to a cycle. Consider a
configuration from which $\theta$ never again removes a frame. Its outgoing path in $W$ is
finite and acyclic, so it ends in a thread $\theta^*$ with no outgoing edge: its top frame
is neither blocked nor waiting on another thread, so it can be reduced. By
\cref{ass:fair} it is reduced, and by termination of method bodies its top frame reaches
a value and is removed by \textsc{Turn-End} or \textsc{Fast-Return}. Induction on the
number of frames above the one that the predecessor of $\theta^*$ waits for shows that
this frame is eventually removed, which unblocks the predecessor; repeating along the
path, $\theta$ progresses --- a contradiction.
\end{proof}

\begin{corollary}[Liveness without \fs{}]
\label[corollary]{cor:asynclive}
If a program never evaluates $\fsend$, every turn completes under \cref{ass:fair}.
\end{corollary}
\begin{proof}
Without \fs{} there are no blocked or waiting frames, so the wait-for graph has no edges.
\end{proof}

This is the liveness half of the isolated turn principle, recovered as the case without
\fs{}. With \fs{}, liveness is a property of the program's pattern of synchronous calls
rather than of the model.

\begin{proposition}[What a per-thread call chain detects]
\label[proposition]{prop:localcycles}
Suppose a thread checks, before each \textsc{Fast-Local} to $\iota_b$, whether $\iota_b$
has a frame on its own stack. The check succeeds exactly for the self-loops
$\theta \to \theta$ of the wait-for graph. It does not detect cycles of length two or
more.
\end{proposition}
\begin{proof}
A self-loop arises when the top frame of $\theta$ is blocked on an actor whose frame is on
$\theta$ itself, which is what the check tests. A cycle through two threads
$\theta_1 \to \theta_2 \to \theta_1$ arises when the top frame of each is blocked on an
actor whose frame is on the other: each check inspects only its own stack, finds the
callee absent, and passes. For example, thread~1 runs a \fs{} to $A$ whose body makes a
\fs{} to $B$, while thread~2 runs a \fs{} to $B$ whose body makes a \fs{} to $A$.
\end{proof}

A per-thread detector, such as the call-chain test designed in
\cite{mayeski2026fastsend},\footnote{\kh{} does not implement the call-chain test yet:
\fs{} takes the receiver's mutex directly, so a synchronous cycle currently deadlocks.}
therefore turns same-thread cycles into errors and leaves cross-thread cycles as
deadlocks. Detecting those needs shared state: a global wait-for graph, a global order on
actors, or bounded waiting with a fallback to asynchronous delivery.

\subsection{Order}

\begin{proposition}[Order]
\label[proposition]{prop:order}
(1) Messages sent with \snd{} from $\iota_a$ to $\iota_b$ start turns of $\iota_b$ in the
order they were sent. (2) A \fs{} from $\iota_a$ to $\iota_b$ starts its turn before any
message already in $\iota_b$'s mailbox, including earlier \snd{}s from $\iota_a$.
(3) Messages sent during the callee's turn of a \fs{} are enqueued before any message the
caller sends after the \fs{} returns.
\end{proposition}
\begin{proof}
(1) \textsc{Send} appends and \textsc{Receive} removes the head. (2) The \fs{} rules start
the callee's turn without consulting its mailbox. (3) The caller's frame is waiting from
the \fs{} until \textsc{Fast-Return}, and the callee's sends happen before that.
\end{proof}

Part (2) is the formal counterpart of ``\fs{} bypasses the mailbox''. On the FPGA runtime
of Section~\ref{sec:substrates} it is implemented by a dedicated input port that each
actor process serves before its mailbox FIFO.

\subsection{\fs{} is not send-then-wait}

Let \emph{send-then-wait} be the pattern of the first group in
Section~\ref{sec:related-sync}: the caller sends a request with \snd{} and its frame then
blocks until a reply message for it arrives. We model the blocking by a frame that no rule
reduces until the reply has been enqueued in the caller's mailbox; the frame stays on its
thread.

\begin{proposition}[Shared resources]
\label[proposition]{prop:notsendwait}
Let actors $\iota_a$ and $\iota_b$ share a home thread, and let $\iota_b$'s method return
without further communication. A turn of $\iota_a$ that obtains $\iota_b$'s result by
send-then-wait never completes. The same turn using \fs{} with \textsc{Fast-Local}
completes.
\end{proposition}
\begin{proof}
With send-then-wait, the request is enqueued in $\iota_b$'s mailbox and $\iota_a$'s frame
stays on the shared thread $\theta$. \textsc{Receive} for $\iota_b$ requires $\theta$ to
have an empty stack, which it never has again, so the reply is never produced. With \fs{},
\textsc{Fast-Local} pushes $\iota_b$'s turn on $\theta$ above $\iota_a$'s frame; it
terminates and \textsc{Fast-Return} resumes $\iota_a$.
\end{proof}

The encodings of synchronous into asynchronous communication
\cite{hondatokoro1991,boudol1992} are sound in a setting where every process progresses
independently. In a setting with a bounded number of execution resources --- actor
groups, ABS cogs, E vats, a pinned thread per core, a fixed number of FPGA processes ---
\fs{} and send-then-wait differ in whether they complete. This is also why systems of the
second group in Section~\ref{sec:related-sync} permit synchronous calls exactly inside a
shared scheduling unit: there, send-then-wait is not an option.

% ---------------------------------------------------------------
\section{Simulation and debugging from a global queue}
\label{sec:debug}
% ---------------------------------------------------------------

A backtest replays recorded inputs through the same actors that trade live. A debugger
needs the same thing plus the ability to stop, step and inspect. Both are easiest when
every actor in the system --- whatever language it is written in and whatever substrate
it normally runs on --- is driven from a single queue by a single thread: every message
goes into the queue, receivers take their messages from it, and the execution is one
deterministic sequence that can be replayed, paused and inspected. The \kh{} actor group,
which co-schedules actors behind one mailbox, already serves as such a simulator for CPU
actors \cite{mayeski2026fastsend}.\footnote{The cross-language, cross-substrate global
queue described in this section is not implemented in \kh{} yet; CPU actors in one group
already run this way.} The question is whether \fs{} breaks the
correspondence between that simulator and the concurrent system.

\subsection{The global-queue semantics}

The global-queue semantics $\mathit{GQ}$ is the calculus of Section~\ref{sec:calculus}
with one thread $\theta_0$, the home thread of every actor, and with per-actor mailboxes
replaced by one FIFO queue $Q$ of pairs $(\iota_b, m)$:
\begin{itemize}[leftmargin=1.4em]
  \item \textsc{Send} appends $(\iota_b, \Msg{msg, \pass(\overline{v})})$ to $Q$;
  \item \textsc{Receive-GQ}: when $\theta_0$ has an empty stack, remove the head
  $(\iota_b, m)$ of $Q$ and start $\iota_b$'s turn for $m$;
  \item \fs{} uses \textsc{Fast-Local} on $\theta_0$.
\end{itemize}
An actor that normally runs on another substrate takes part through its method bodies
alone: an FPGA actor's turn is its step function, executed by the simulator in software;
a Rust actor's turn is its handler, called across the language boundary. By receiver
transparency, the body is the same whatever starts the turn.

The production semantics $\mathit{P}$ is the calculus of Section~\ref{sec:calculus} with
any assignment of actors to threads and any choice of \textsc{Fast-Local} or
\textsc{Fast-Remote} per call site, subject to one condition that the runtimes studied
here satisfy:

\begin{assumption}[Dedicated remote resources]
\label[assumption]{ass:dedicated}
A thread that serves as the home thread of a \textsc{Fast-Remote} callee is the home of
that actor only.
\end{assumption}

On the FPGA runtime each actor is its own process, so the assumption holds for \fs{}
calls into the FPGA. CPU actors that share a thread are called with \textsc{Fast-Local}.

\subsection{Soundness}

Say that two executions \emph{agree} if every actor starts the same turns in the same order,
with $\alpha$-equivalent messages, and every turn ends with an $\alpha$-equivalent value
and heap.

\begin{theorem}[Global-queue simulation is sound]
\label{thm:gqsound}
Under \cref{ass:dedicated}, for every finite execution of $\mathit{GQ}$ from a given
initial configuration there is an execution of $\mathit{P}$ from the corresponding initial
configuration that agrees with it. In particular, a configuration in which a turn is
stuck, or an error state, reached in $\mathit{GQ}$ has an agreeing counterpart in
$\mathit{P}$.
\end{theorem}
\begin{proof}
We build the $\mathit{P}$ execution step by step, maintaining the invariant that at every
point where $\theta_0$ has an empty stack in $\mathit{GQ}$, every thread of $\mathit{P}$
has an empty stack, the heaps agree, and each actor's mailbox in $\mathit{P}$ is the
projection of $Q$ onto that actor, in order.

\emph{Receive.} $\mathit{GQ}$ takes the head $(\iota_b, m)$ of $Q$. By the invariant,
$m$ is the head of $\iota_b$'s mailbox in $\mathit{P}$, since every earlier pair for
$\iota_b$ in $Q$ has already been removed. $\iota_b$'s home thread is idle and $\iota_b$
is not busy, so \textsc{Receive} applies in $\mathit{P}$ to the same message.

\emph{Within a turn.} $\mathit{P}$ reduces only the thread that holds the corresponding
frame, mirroring each $\mathit{GQ}$ step: actor-local steps are identical; \textsc{Send}
appends to the receiver's mailbox, which preserves the projection invariant.

\emph{\fs{}.} In $\mathit{GQ}$, \textsc{Fast-Local} to $\iota_b$ applies if and only if
$\iota_b$ has no frame on $\theta_0$'s stack. In the mirrored $\mathit{P}$ execution the
only frames in existence are the ones mirroring that stack, so $\iota_b$ is busy in
$\mathit{P}$ exactly when it is busy in $\mathit{GQ}$. If $\mathit{P}$ uses
\textsc{Fast-Local} for the call, it applies. If it uses \textsc{Fast-Remote}, the rule
also needs $\iota_b$'s home thread to be idle. By \cref{ass:dedicated} a frame reaches
that thread only as a turn of $\iota_b$ or stacked above one; $\iota_b$ has no frame, so
the thread is idle. The callee's turn then
proceeds as above, and \textsc{Fast-Return} applies in both.

If $\mathit{GQ}$ is stuck on a \fs{} to a busy actor, the mirrored $\mathit{P}$ frame is
blocked on the same actor. Agreement of values and heaps follows from
\cref{thm:determinism}, since the mirrored turns receive the same replies.
\end{proof}

\begin{corollary}[What the simulator shows is real]
Every value a \fs{} returns in $\mathit{GQ}$, every assertion failure, error state or
same-thread \fs{} cycle it reaches, can occur in the concurrent system. A $\mathit{GQ}$
run is deterministic given its inputs, so it can be replayed exactly.
\end{corollary}

So \fs{} does not break the global-queue simulation. What makes the argument go through is
that in a single-threaded run the callee of a \fs{} is idle unless it is already on the
current call chain: the simulator only ever takes the branch of \fs{} in which the callee
was free, which is one of the branches the concurrent system can take. It is not obvious
in general: without \cref{ass:dedicated}, a \textsc{Fast-Remote} callee whose resource is
occupied by another actor blocks in $\mathit{P}$ where $\mathit{GQ}$ proceeds, and the
simulator would then show behaviour the concurrent system cannot produce.

\subsection{What the simulator cannot show}

\begin{proposition}[Incompleteness]
\label[proposition]{prop:gqincomplete}
There are executions of $\mathit{P}$ that agree with no execution of $\mathit{GQ}$. In
particular: (1) a deadlock on a wait-for cycle through two or more threads; (2) a turn of
$\iota_a$ that makes a \fs{} to $\iota_b$ after another thread has run a turn of
$\iota_b$ that started after $\iota_a$'s turn began.
\end{proposition}
\begin{proof}
In $\mathit{GQ}$ only one thread exists, so the wait-for graph can contain only
self-loops, and no turn starts while another is in progress except as a \fs{} callee on
the same stack. Both executions described require a turn to start on another thread while
$\iota_a$'s turn is in progress.
\end{proof}

A backtest can therefore confirm bugs but cannot rule out the two classes above. The
first can be approached offline:

\begin{proposition}[Cycles in the call graph]
\label[proposition]{prop:callgraph}
Let $G$ be the directed graph on actors with an edge $\iota_a \to \iota_b$ whenever some
turn of $\iota_a$ can evaluate a \fs{} to $\iota_b$. If $G$ is acyclic, no execution of
$\mathit{P}$ reaches a wait-for cycle.
\end{proposition}
\begin{proof}
The frames on any thread, bottom to top, belong to a path in $G$, since each frame above
the bottom one was pushed by a \fs{} from the frame below (\textsc{Fast-Local}); a
\textsc{Fast-Remote} callee's frame extends the same path onto another thread. A wait-for
edge $\theta \to \theta'$ adds the edge from the actor on top of $\theta$ to an actor
already on the path held by $\theta'$. Following a wait-for cycle and concatenating these
paths yields a cycle in $G$.
\end{proof}

A $\mathit{GQ}$ run can record the \fs{} edges it executes. The recorded graph is a
subgraph of $G$ --- it contains only the calls the inputs exercised --- so a cycle in it is
a real risk to be removed, and the absence of a cycle is evidence only for the paths the
backtest covered. A static analysis of the call sites would give all of $G$.

\paragraph{Time.} $\mathit{GQ}$ reproduces what happens, not when. Latency on each
substrate is a separate question, answered by measurement or by the substrate's own
timing model (Section~\ref{sec:eval}).

% ---------------------------------------------------------------
\section{Mapping \snd{} and \fs{} onto substrates}
\label{sec:substrates}
% ---------------------------------------------------------------

The calculus leaves open what a thread is and how a mailbox is built. Each substrate
supplies both, and each supplies a distinct mechanism for each primitive
(\cref{tab:mapping}).

\begin{table}[t]
\centering
\caption{\snd{} and \fs{} on each substrate and between substrates, and the rule of
Section~\ref{sec:calculus} each \fs{} implements.}
\label{tab:mapping}
\small
\begin{tabular}{@{}p{3.2cm}p{4.6cm}p{6.3cm}@{}}
\toprule
Where & \snd{} & \fs{} \\
\midrule
CPU, one process & message into the receiver's mailbox; a thread per actor or per group & callee's handler runs on the caller's thread under the callee's lock (\textsc{Fast-Local}) \\
C++ $\leftrightarrow$ Rust, one process & message converted and enqueued through a C ABI bridge & callee's handler runs on the caller's thread through the bridge (\textsc{Fast-Local}) \\
FPGA, one process & --- & callee is an actor object inside the caller's process; its handler is inlined into the caller's pipeline (\textsc{Fast-Local}) \\
FPGA, between processes & on-chip router into the receiver's mailbox FIFO & request and reply relayed through the host (\textsc{Fast-Remote}) \\
CPU $\to$ FPGA & envelope over PCIe, router, mailbox FIFO & envelope to the actor's fast\_send port, served before its mailbox; the CPU thread waits for the reply (\textsc{Fast-Remote}) \\
FPGA $\to$ CPU & envelope over PCIe; a host bridge thread enqueues it & request over PCIe; the bridge thread runs the CPU handler inline; the FPGA process stalls until the reply returns (\textsc{Fast-Remote}) \\
GPU (design) & ring buffer in device memory & device function call inside one kernel (\textsc{Fast-Local}); host calls wait on a reply ring (\textsc{Fast-Remote}) \\
Another machine & network message & request/reply RPC; the caller waits (\textsc{Fast-Remote}) \\
\bottomrule
\end{tabular}
\end{table}

The pattern is the same everywhere. \snd{} is a queue between two independently scheduled
resources. \fs{} is either an inline call on the caller's resource, where the substrate
lets one resource run another actor's handler, or a blocking request to a resource
dedicated to the callee, where it does not. The calculus of Section~\ref{sec:calculus}
covers both with one semantics, and \cref{thm:isolation,thm:determinism,thm:liveness}
hold for any mixture.

A \fs{} is available between any two substrates where the caller can wait. A CPU thread
can wait on a condition. An FPGA process can wait on a blocking stream read, which stalls
that process alone while the rest of the chip keeps running. What differs between
substrate pairs is the cost of the wait, which Section~\ref{sec:eval} measures where it can.

% ---------------------------------------------------------------
\section{Implementations}
\label{sec:impl}
% ---------------------------------------------------------------

\paragraph{CPU runtime.} The \kh{} C++ runtime \cite{mayeski2026fastsend} implements
actors with a mutex-protected turn, a selectable mailbox per actor, \snd{}, \fs{} and
groups. \fs{} takes the callee's lock on the caller's thread and runs its handler
inline; a reply is returned as a value. A delivery flag is private to the framework, and
the two queue fields a handler may read are defined identically for both delivery modes,
so the handler receives no information about how it was reached. A mailbox may also be
configured to busy-poll instead of sleeping, which changes the cost of \snd{} between
threads but not its semantics.

\paragraph{Rust runtime and interop.} A separate Rust actor runtime hosts Rust actors.
C++ and Rust actors reach each other through generated C-ABI bridges, from one message
schema; both \snd{} and \fs{} work across the boundary.\footnote{\kh{} does not return
\fs{} replies across the language boundary yet; a cross-language \fs{} delivers the
message and returns no value.}

\paragraph{FPGA runtime.} The FPGA runtime is written in Vitis HLS C++. Each actor is a
hardware process with a mailbox FIFO, a \fs{} port served before the mailbox, and a
port on which replies to its own remote \fs{}s arrive; its handlers are selected by a
switch on the message identifier. A router delivers messages by destination actor, using a
route table fixed when the design is built, and a host link connects the design to PCIe.
On the CPU, an \texttt{ActorRef} can refer to an FPGA actor; \snd{} and \fs{} on it write
envelopes to the card, and a bridge thread delivers messages and \fs{} requests from the
card to CPU actors, passing a stand-in for the FPGA actor as the sender so that ordinary
replies return to the chip. The design runs today in software, with one thread per
hardware process.\footnote{The FPGA runtime has not yet been synthesized or run on a
card; \fs{} between two FPGA processes is relayed through the host.}

\paragraph{GPU runtime.} The same structure maps onto a GPU: actors as warps or thread
blocks of a persistent kernel, mailboxes as device-memory rings, \fs{} inside the kernel as
a device function call, and a host bridge over rings in pinned memory. It is a design; we
have not built it.

% ---------------------------------------------------------------
\section{Evaluation}
\label{sec:eval}
% ---------------------------------------------------------------

\paragraph{CPU.} \Cref{tab:cpu} gives the round-trip cost of a ping/pong exchange with a
reply on an AMD EPYC 9374F (RHEL 9.2, \texttt{g++ -O3 -march=native}, two dedicated
cores). The four rows are four implementations of the same exchange in the calculus:
\snd{} between two threads that sleep when idle, the same with threads that busy-poll,
\snd{} within one group, and \fs{}.

\begin{table}[t]
\centering
\caption{Round trip of one ping/pong with a reply, median over runs, AMD EPYC 9374F. The
\fs{} figure is amortized over $10^6$ round trips because its per-sample value is below
the clock's resolution of about 10~ns.}
\label{tab:cpu}
\small
\begin{tabular}{@{}lrl@{}}
\toprule
Implementation & ns & What it pays for \\
\midrule
\snd{}, two threads, sleeping & $\sim\!3370$ & two cross-core handoffs and two thread wake-ups \\
\snd{}, two threads, busy-polling & $370$ & two cross-core handoffs; each thread holds a core \\
\snd{}, one group (one thread) & $90$ & two queue operations \\
\fs{} & $\sim\!30$ & two inline calls under the callee's lock \\
\bottomrule
\end{tabular}
\end{table}

Most of the cost of \snd{} between threads is putting a thread to sleep and waking it; the
cross-core handoff itself is about 185~ns each way. \fs{} costs about one mispredicted
virtual call per hop (9.1~ns against 7.5--9.3~ns, measured the same way) and about
$5\times$ a direct call \cite{mayeski2026fastsend}.

\paragraph{Across languages.} On an Apple Silicon machine, a one-way \fs{} without a reply
costs about 13--15~ns within one language, 46~ns from C++ to Rust and 70~ns from Rust to
C++; the difference is the marshalling, the C-ABI call, and a per-call lookup of the
receiver by name.

\paragraph{FPGA.} The FPGA runtime's semantics are tested in software: \snd{} through the
router, \fs{} inside a process, and \snd{} and \fs{} in both directions between CPU and
FPGA actors, with ordering, error reporting and the fast\_send port's priority over the
mailbox checked by unit tests. Its latencies are not yet measured. The per-part cycle
counts will come from high-level synthesis, and the PCIe component from measurement on a
card; we do not report either here.

% ---------------------------------------------------------------
\section{Related work}
\label{sec:related}
% ---------------------------------------------------------------

\paragraph{Actor semantics.} The operational semantics of Agha, Mason, Smith and Talcott
\cite{agha1997} and the four-family semantics of De Koster and De Meuter
\cite{dekoster2024} contain only asynchronous communication. We extend the latter.
Our isolation and liveness results are stated against their isolated turn principle.

\paragraph{Synchronous calls.} ABCL/1 \cite{yonezawa1986} and ABS \cite{johnsen2012abs}
provide synchronous calls between concurrent objects that block the caller while the
callee runs on its own processor, and ABS adds reentrant synchronous calls inside a cog.
E \cite{miller2005} restricts immediate calls to one vat. Extended Rebeca
\cite{sirjani2005extended} adds synchronous message passing to the Rebeca actor language
with a formal semantics; comparing its scheduling of synchronous messages with \fs{} is
left for future work. The survey of active-object languages by De Boer et al.
\cite{deboer2017survey} maps this design space.

\paragraph{Encodings.} Honda and Tokoro \cite{hondatokoro1991} and Boudol
\cite{boudol1992} encode synchronous in asynchronous communication; \cref{prop:notsendwait}
locates where such encodings stop applying once resources are bounded.

\paragraph{Actors on hardware.} CAL and the StreamBlocks compiler
\cite{bezati2021streamblocks} compile dataflow actors to both CPUs and FPGAs. Dataflow
actors fire on token availability and have no synchronous call; our FPGA runtime targets
the message-passing actor model with both primitives.

% ---------------------------------------------------------------
\section{Limitations and future work}
\label{sec:future}
% ---------------------------------------------------------------

\begin{itemize}[leftmargin=1.4em]
  \item The calculus and proofs are on paper. A mechanisation, for instance by extending
  the PLT Redex models of \cite{dekoster2024}, would check them.
  \item The calculus has no failures; a \fs{} to an actor on a failed resource waits
  forever. Timeouts and their interaction with isolation are open.
  \item Detecting cross-thread cycles at run time needs shared state
  (\cref{prop:localcycles}); its cost on the hot path is not measured.
  \item FPGA latencies await synthesis and a card; the GPU runtime is a design.
  \item Extended Rebeca's synchronous messages should be compared with \fs{} in detail.
\end{itemize}

% ---------------------------------------------------------------
\section{Conclusion}
% ---------------------------------------------------------------

\fs{} adds a second primitive to the actor model: the receiver's turn runs now, on a
resource the caller waits on, and returns a value. Formally, it keeps the model's isolation,
makes determinism relative to replies, and replaces unconditional liveness with liveness
exactly when waits are acyclic; asynchronous send alone is the acyclic case. It is not
reducible to send-then-wait once execution resources are shared. A single global queue
remains a sound simulator with \fs{}, so a backtest or a debugger that runs every actor,
in any language and on any substrate, from one queue shows only behaviour the real system
can produce, though not all of it. And on every substrate we have built, \snd{} and \fs{}
are two different mechanisms --- a queue, and an inline or blocking call --- which is the
practical case for treating both as primitives of the model.

% ---------------------------------------------------------------
\begin{thebibliography}{99}
\bibitem{agha1986} G.~Agha. \emph{Actors: A Model of Concurrent Computation in Distributed
  Systems.} MIT Press, 1986.
\bibitem{agha1997} G.~A. Agha, I.~A. Mason, S.~F. Smith, and C.~L. Talcott. A foundation
  for actor computation. \emph{Journal of Functional Programming}, 7(1):1--72, 1997.
\bibitem{bezati2021streamblocks} E.~Bezati et al. StreamBlocks: A compiler for
  heterogeneous dataflow computing. Technical report, EPFL, 2021.
\bibitem{boudol1992} G.~Boudol. Asynchrony and the $\pi$-calculus. Technical Report 1702,
  INRIA Sophia-Antipolis, 1992.
\bibitem{deboer2017survey} F.~de Boer, V.~Serbanescu, R.~H\"ahnle, L.~Henrio, J.~Rochas,
  C.~C. Din, E.~B. Johnsen, M.~Sirjani, E.~Khamespanah, K.~Fernandez-Reyes, and A.~M.
  Yang. A survey of active object languages. \emph{ACM Computing Surveys}, 50(5), 2017.
\bibitem{dekoster2024} J.~De Koster and W.~De Meuter. A formal specification for half a
  century of actor systems. Manuscript, Vrije Universiteit Brussel, 2023.
\bibitem{hewitt1973} C.~Hewitt, P.~Bishop, and R.~Steiger. A universal modular ACTOR
  formalism for artificial intelligence. In \emph{IJCAI}, 1973.
\bibitem{hondatokoro1991} K.~Honda and M.~Tokoro. An object calculus for asynchronous
  communication. In \emph{ECOOP '91}, LNCS 512, pages 133--147, 1991.
\bibitem{johnsen2012abs} E.~B. Johnsen, R.~H\"ahnle, J.~Sch\"afer, R.~Schlatte, and
  M.~Steffen. ABS: A core language for abstract behavioral specification. In \emph{Formal
  Methods for Components and Objects (FMCO 2010)}, LNCS 6957, pages 142--164, 2012.
\bibitem{mayeski2026fastsend} V.~Mayeski. Evaluating suitability of the actor model of
  concurrency for high-frequency trading: the \texttt{fast\_send} mechanism and a
  tick-to-book latency study. arXiv:2609.21173, 2026.
\bibitem{miller2005} M.~S. Miller, E.~D. Tribble, and J.~Shapiro. Concurrency among
  strangers: programming in E as plan coordination. In \emph{Trustworthy Global
  Computing}, 2005.
\bibitem{sirjani2005extended} M.~Sirjani, F.~S. de Boer, A.~Movaghar, and A.~Shali.
  Extended Rebeca: a component-based actor language with synchronous message passing. In
  \emph{Application of Concurrency to System Design (ACSD)}, 2005.
\bibitem{yonezawa1986} A.~Yonezawa, J.-P. Briot, and E.~Shibayama. Object-oriented
  concurrent programming in ABCL/1. In \emph{OOPSLA '86}, pages 258--268, 1986.
\end{thebibliography}

\end{document}
